\section{Infinitesimal lifting and section growth}\label{bl:sec:lifting}

We now apply the residue insertion and the two frame assertions to
prove lifting through every finite order. Finite-cover invariants
and the weight-zero part of the frame action will then return the
lifted layers to \(V\), where their positive twists force unbounded
section spaces.

Retain the complex algebraic setup of the preceding sections.
In particular, \(I=\OO_Z(-E)=y\OO_Z\), the pullback weight of \(y\) is \(-1\), and
\(\Omega=\res(\sigma)\) has weight \(a=r/q>0\).  Write
\begin{equation}\label{bl:lift:insertion}
 \iota:R^1g_*\OO_E\hookrightarrow
 R^1h_*\omega_R(D_h)=F_0\mathcal N\hookrightarrow\mathcal N
\end{equation}
for the injections of Lemmas~\ref{bl:lem:residue-injection}
and~\ref{bl:lem:lowest-piece}.  Thus \(\iota(e)=\Omega e\), with pullback and
residue understood.  If \(F\in g_*\OO_E\), the expression \(\Omega Fe\)
means \(\iota(Fe)\): multiplication takes place on \(E\), before the
insertion.  This convention does not require \(F\) to come from \(B\).

For all integers \(j<d\), we push the underlying sheaf of complex vector
spaces of \(I^j/I^d\) along \(|g|\) using the Zariski-topology convention
of Section~\ref{bl:cov:obstruction-contract}; negative powers mean
invertible fractional ideals. No morphism from an infinitesimal
thickening to \(B\) is assumed. In contrast, multiplication by \(y^j\) identifies
\(I^j/I^{j+1}\) with the coherent sheaf \(\OO_E\), with an equivariant
weight shift.  Its additive higher direct images are therefore the
underlying sheaves of the usual coherent higher direct images under the
projective morphism \(g\).

\begin{proposition}[Infinitesimal lifting]\label{bl:prop:lifting}
For every \(k\geq1\), the truncation map
\begin{equation}\label{bl:lift:successive}
 g_*(\OO_Z/I^{k+1})\longrightarrow g_*(\OO_Z/I^k)
\end{equation}
is surjective as a map of additive sheaves on \(B\).  More generally, for
all integers \(j<d\), the maps
\[
 g_*(I^j/I^{d+1})\longrightarrow g_*(I^j/I^d),
 \qquad
 g_*(I^j/I^d)\longrightarrow g_*(I^j/I^{j+1})
\]
are surjective.  Each of these surjections is also surjective on sections
over every affine open subset of \(B\).
\end{proposition}

\begin{proof}
\emph{The obstruction derivation.}
We first prove \eqref{bl:lift:successive} by induction on \(k\).
Put \(A_d=g_*(\OO_Z/I^d)\).  The obstruction to surjectivity at the
\(k\)-th step is the connecting homomorphism
\[
 \partial_k:A_k\longrightarrow R^1g_*(I^k/I^{k+1}).
\]
The previous steps give a surjection \(A_k\to A_1\).  For \(k>1\), its
kernel is \(g_*(I/I^k)\), by left exactness.  This kernel lifts to
\(g_*(I/I^{k+1})\): multiplication by \(y\) identifies the required
surjection with the already established map \(A_k\to A_{k-1}\).
Consequently \(\partial_k\) vanishes on that kernel.  For \(k=1\) the
kernel is zero.  After framing the target layer by \(y^k\), we obtain in
either case a uniquely defined map
\begin{equation}\label{bl:lift:derivation}
 \delta:g_*\OO_E\longrightarrow R^1g_*\OO_E,
 \qquad
 \partial_k(\widetilde F)=y^k\delta(F).
\end{equation}
Here \(\widetilde F\) denotes any local lift of \(F\) to \(A_k\);
we suppress the index \(k\) in the notation \(\delta_k\) of
Section~\ref{bl:cov:obstruction-contract}.

The map \(\delta\) is a \(\C\)-linear derivation, where
\(R^1g_*\OO_E\) is a module over \(g_*\OO_E\) by multiplication on
\(E\).  To check this assertion, work locally on the base, choose lifts
of \(F,G\) to \(A_k\), and represent them on an open cover along \(E\)
by ambient functions \(F_\alpha,G_\alpha\).  We use the \v{C}ech
convention \((dc)_{\alpha\beta}=c_\beta-c_\alpha\).
Both differences \(F_\beta-F_\alpha\) and
\(G_\beta-G_\alpha\) are divisible by \(y^k\).  In
\[
 F_\beta G_\beta-F_\alpha G_\alpha
 =F_\alpha(G_\beta-G_\alpha)
  +G_\alpha(F_\beta-F_\alpha)
  +(F_\beta-F_\alpha)(G_\beta-G_\alpha),
\]
the last term vanishes modulo \(y^{k+1}\), since \(2k\geq k+1\).
Divide by \(y^k\) and restrict to \(E\).  The resulting cocycles give
\begin{equation}\label{bl:lift:leibniz}
 \delta(FG)=F\delta(G)+G\delta(F),\qquad \delta(1)=0.
\end{equation}
The factorization just proved ensures independence of the chosen shorter
lifts.

\smallskip\noindent
\emph{The graph residue identity.}
We relate this derivation to right differentiation on \(\mathcal N\).
Choose a coordinate chart in \(B\), with coordinates
\(t_1,\ldots,t_b\), and put \(e_i=\delta(t_i)\), applying \(\delta\)
to their pullbacks to \(E\).  For every local section
\(F\in g_*\OO_E\), we claim that
\begin{equation}\label{bl:lift:variable-identity}
 \iota\bigl(\delta(F)\bigr)
   +\sum_{i=1}^b\iota(Fe_i)\,\partial_{t_i}=0
 \qquad\text{in }\mathcal N.
\end{equation}
It is essential here that \(F\) is arbitrary.

This can be checked at base germs.  Shrink the base so that the induction
provides lifts of \(F,t_1,\ldots,t_b\) modulo \(I^k\), and choose
ambient representatives \(F_\alpha,v_{\alpha i}\) on an open cover
along \(E\).  Pull them back to \(Y\), still writing \(y\) for its
pullback.  On overlaps write
\begin{equation}\label{bl:lift:overlap-functions}
 v_{\beta i}-v_{\alpha i}=y^k u_{\alpha\beta,i},
 \qquad F_\beta-F_\alpha=y^k b_{\alpha\beta}.
\end{equation}
The restrictions of \(u_{\alpha\beta,i}\) and
\(b_{\alpha\beta}\) to \(E\) represent \(e_i\) and \(\delta(F)\).
The same assertions hold if representatives are specified only modulo
\(y^{k+1}\).  On the product with the coordinate chart of \(B\), put
\(\psi_{\alpha i}=t_i-v_{\alpha i}\) and
\(dt=dt_1\wedge\cdots\wedge dt_b\).  In the full local-cohomology
module \(\mathcal M\) of Lemma~\ref{bl:lem:lowest-piece}, consider
the zero-cochain of ordered generalized fractions
\begin{equation}\label{bl:lift:cochain}
 C_\alpha=
 \left[\frac{\sigma F_\alpha\wedge dt}
 {y^k\prod_i\psi_{\alpha i}}\right].
\end{equation}
The brackets include the divisor support direction as well as the
\(b\) graph directions, in a fixed order.  Poles along \(D_h\) are
allowed.  The divisor pole bound for \eqref{bl:lift:cochain} is
\(R+D_h+k\widehat E\), so these classes belong to \(\mathcal M\);
they need not belong to \(F_0\mathcal M\).

Here is a calculation valid also at crossings and at components with
multiplicity in \(\widehat E\).  Locally invert an equation of \(D_h\)
and let \(r_0\) be a reduced equation for \(R\).  In this localization,
\(\sigma=\alpha/r_0\), where \(\alpha\) is a regular top form,
\(r_0\mid y\), and the zero support of \(y\) is exactly \(R\).
The divisor denominator in \eqref{bl:lift:cochain} is \(r_0y^k\).
In the quotient by this denominator,
\begin{equation}\label{bl:lift:nilpotent-modulus}
 y^{2k}=0,\qquad y^{k+1}=0
 \quad\text{in }\OO_Y(*D_h)/(r_0y^k),
\end{equation}
because \(r_0\mid y\) and \(k\geq1\).  Thus the ideal generated by
\(y^k\) is square zero there.  Localizations by
\(\psi_{\alpha i}\) and \(\psi_{\beta i}\) agree canonically,
since the two elements differ by a nilpotent element, and the exact
first-order formula is
\[
 \psi_{\beta i}^{-1}
 =\psi_{\alpha i}^{-1}
   +y^k u_{\alpha\beta,i}\psi_{\alpha i}^{-2}.
\]
Multiplying these identities and using
\(F_\beta=F_\alpha+y^kb_{\alpha\beta}\), all terms containing two
differences vanish by \eqref{bl:lift:nilpotent-modulus}.  After dividing
by the divisor denominator, the result in local cohomology is
\begin{align}\label{bl:lift:cech-expansion}
 C_\beta-C_\alpha
 ={}&\left[
  \frac{\sigma b_{\alpha\beta}\wedge dt}
       {\prod_j\psi_{\alpha j}}\right] \notag\\
 &+\sum_i\left[
  \frac{\sigma F_\alpha u_{\alpha\beta,i}\wedge dt}
       {\psi_{\alpha i}^{2}\prod_{j\ne i}\psi_{\alpha j}}
                    \right].
\end{align}
One can formalize this passage by first computing graph local cohomology
over \(\OO_Y(*D_h)/(r_0y^k)\), and then mapping to divisor local
cohomology by the fraction with denominator \(r_0y^k\).
Changing the graph equations by the indicated nilpotents gives the same
localizations at the first stage.  Divisor local cohomology is the direct
limit of these denominator calculations, so the resulting identity is
an identity in the original support module.  This also explains why
ordinary reduced-support division, without the factor \(y^k\) in the
modulus, would not justify the computation.

The coefficients remaining in \eqref{bl:lift:cech-expansion} have only the
simple divisor pole bound \(R+D_h\).  Their residues are respectively
\(\Omega\delta(F)\) and \(\Omega Fe_i\) as \v{C}ech classes under
\eqref{bl:lift:insertion}.  Indeed division by \(y^k\) has already
cancelled in the displayed coefficients.  Changing a coefficient
representative by a multiple of \(y\) changes its product with
\(\sigma\) by a form regular in the \(R\)-direction after the
localization, because \(r_0\mid y\).  Its residue is therefore zero.
This is the product-residue map on the entire divisor, including its
crossings.

For completeness, the sign in \eqref{bl:lift:cech-expansion} agrees with
the right-module action.  If \(\eta\) is a form pulled back from
\(Y\), right differentiation of top forms is minus differentiation of
their coefficients.  Hence
\begin{equation}\label{bl:lift:right-sign}
 \left[\frac{\eta\wedge dt}{\prod_j\psi_{\alpha j}}\right]
       \partial_{t_i}
 =\left[\frac{\eta\wedge dt}
 {\psi_{\alpha i}^{2}\prod_{j\ne i}\psi_{\alpha j}}\right].
\end{equation}
The rightmost-term map from \(\mathcal M[0]\) to the relative right
de Rham complex is a map of complexes and commutes with target
differentiation.  Compute its derived pushforward with an affine
\v{C}ech refinement.  The total differential of the zero-cochain
\eqref{bl:lift:cochain}, which lies in relative degree zero, is its
\v{C}ech coboundary, since there is no subsequent relative term.
Its image in \(\mathcal N=\mathcal H^1\operatorname{pr}_{B,+}
\mathcal M\) is zero.  Equations~\eqref{bl:lift:cech-expansion}
and~\eqref{bl:lift:right-sign}, and the lowest-piece identification in
Lemma~\ref{bl:lem:lowest-piece}, prove \eqref{bl:lift:variable-identity}.
No ambient lift of the map \(g\) on a neighborhood of a whole fiber
was used: the chosen ambient representatives and their overlaps suffice.

\smallskip\noindent
\emph{Vanishing of the obstruction.}
Apply \eqref{bl:lift:variable-identity} first with \(F=1\), and take the
first symbol. The tensor \(v_k\), written \(v\) at this fixed order, satisfies
\begin{equation}\label{bl:lift:symbol-vector}
 v=\sum_i\iota(e_i)\otimes\partial_{t_i}
 \in H^0(B,F_0\mathcal N\otimes T_B),
 \qquad \operatorname{symbol}(v)=0.
\end{equation}
The restriction of \(\delta\) to \(\OO_B\) factors through the
algebraic K\"ahler differentials \(\Omega_B^1\), by
\eqref{bl:lift:leibniz}.  Thus \eqref{bl:lift:symbol-vector} is independent
of coordinates and is a global algebraic tensor.
The connecting map before the framing in \eqref{bl:lift:derivation} is
equivariant.  Since \(\lambda^*y=\lambda^{-1}y\), the framed map
satisfies
\[
 \lambda^*\delta(F)=\lambda^k\delta(\lambda^*F).
\]
Consequently the derivation tensor has weight \(k\): a coordinate's
weight in its coefficient is cancelled by the weight of its tangent
vector.  Multiplication by \(\Omega\) adds \(a\), so \(v\) has
weight
\begin{equation}\label{bl:lift:positive-weight}
 m=a+k>0.
\end{equation}

Proposition~\ref{bl:prop:frame-positivity} forces the horizontal projection
of \(v\) to vanish.  The vertical tangent line of \(B\to T\) is
freely generated by the invariant, nowhere-zero Euler field
\(\varepsilon=r\beta\partial_\beta\).  Thus
\(v=s_0\otimes\varepsilon\) for a global section
\(s_0\in H^0(B,F_0\mathcal N)\) of weight \(m\).
Write \(\varepsilon=\sum_i c_i\partial_{t_i}\).  The identity with
\(F=1\), before taking the symbol, reads
\[
 0=\sum_i\iota(e_i)\partial_{t_i}
   =\sum_i(s_0c_i)\partial_{t_i}=s_0\varepsilon.
\]
Coefficients occur before the derivatives, so the last equality is
associativity of the right action; no coefficient is commuted through a
derivative.  The Euler assertion of
Proposition~\ref{bl:prop:frame-positivity} now gives
\(0=s_0\varepsilon=-ms_0\), whence \(s_0=0\).
It follows that every \(\iota(e_i)\), and hence every \(e_i\), is
zero.  Return to \eqref{bl:lift:variable-identity} for arbitrary \(F\).
It yields \(\iota(\delta(F))=0\), and injectivity of
\eqref{bl:lift:insertion} gives \(\delta(F)=0\).  This proves the
induction step and hence \eqref{bl:lift:successive}.  The same argument
includes \(\dim T=0\), when there is no horizontal tangent direction.

\smallskip\noindent
\emph{Signed powers and affine sections.}
For any integers \(j<d\), multiplication by the global meromorphic
function \(y^j\) identifies
\[
 \OO_Z/I^{d-j}\simeq I^j/I^d.
\]
It transports the established surjections to all successive Laurent
truncations.  Iteration gives the surjection to the leading layer.
Moreover, each single successive truncation fits into an exact sequence
of additive sheaves
\begin{equation}\label{bl:lift:coherent-kernel}
 0\longrightarrow g_*(I^d/I^{d+1})
 \longrightarrow g_*(I^j/I^{d+1})
 \longrightarrow g_*(I^j/I^d)\longrightarrow0.
\end{equation}
The kernel is coherent on \(B\).  On an affine open it has vanishing
first cohomology, so the additive-sheaf cohomology sequence gives
surjectivity on sections.  Finitely many such steps give the asserted
affine-section surjectivity to every leading layer.  Only the kernels
in \eqref{bl:lift:coherent-kernel} were assigned coherent \(\OO_B\)-module
structures.
\end{proof}

\begin{lemma}[Descent and section growth]\label{bl:lem:descent-growth}
For \(j\in\Z\), define the divisorial sheaf and its layer by
\[
 J_j^V=\OO_V\bigl(-\lceil jG/r\rceil\bigr),
 \qquad \mathcal B_j=J_j^V/J_{j+1}^V,
 \qquad Q_j=f_*\mathcal B_j.
\]
The sheaves \(\mathcal B_j\) are coherent on \(S\), and for every
\(d>j\) the map
\begin{equation}\label{bl:lift:descent-surjection}
 f_*(J_j^V/J_d^V)\longrightarrow Q_j
\end{equation}
is surjective as a map of additive sheaves on \(T\).  Here the left
side is pushed along \(|f|\) on its topological support \(|S|\).
Furthermore,
\[
 h^0\bigl(V,\OO_V(NG)/\OO_V\bigr)\longrightarrow\infty
 \qquad\text{as }N\longrightarrow\infty.
\]
\end{lemma}

\begin{proof}
\emph{The invariant layers.}
Since \(r\) is divisible by every \(d_i\), we have \(0<d_i/r\leq1\).
Thus the coefficientwise difference
\(\lceil(j+1)G/r\rceil-\lceil jG/r\rceil\) has coefficients zero
or one.  A function in the reduced ideal of \(S\) has order at least
one along each \(S_i\).  The divisorial inequalities therefore show
that this ideal carries \(J_j^V\) into \(J_{j+1}^V\).
Consequently \(\mathcal B_j\) is a coherent \(\OO_S\)-module.

Let \(\mu_r\) denote the deck group of the first, full root cover
\(\pi:Z\to P\).  We have
\begin{equation}\label{bl:lift:invariant-ideals}
 (\pi_*I^j)^{\mu_r}=p^*J_j^V
 \qquad (j\in\Z).
\end{equation}
Indeed an invariant meromorphic function in the full root algebra
descends to \(P\), even when that algebra is disconnected.
At a prime over \(p^{-1}(S_i)\), its downstairs order \(n\) becomes
\((r/d_i)n\), whereas \(y\) has order one.  Membership in \(I^j\)
is thus equivalent there to
\[
 (r/d_i)n\geq j
 \quad\Longleftrightarrow\quad n\geq\lceil jd_i/r\rceil.
\]
At other primes it means ordinary regularity.  The codimension-one
test on the normal schemes proves \eqref{bl:lift:invariant-ideals}, also
for negative \(j\).  The identification with the pullback divisorial
sheaf can be checked on the local product charts of \(P\to V\).
Finite pushforward is exact, and averaging over \(\mu_r\) is exact
in characteristic zero.  Hence
\begin{equation}\label{bl:lift:invariant-quotients}
 \bigl(\pi_*(I^j/I^d)\bigr)^{\mu_r}
   =p^*(J_j^V/J_d^V) \qquad(j<d).
\end{equation}

\smallskip\noindent
\emph{Descending the lifting maps.}
To prove \eqref{bl:lift:descent-surjection}, let \(U\subset T\) be
affine and put \(B_U=B\times_T U\).  The morphism \(B\to T\) is
affine, so \(B_U\) is affine.  Define the open neighborhood
\[
 V_U=V\setminus\bigl(S\setminus f^{-1}(U)\bigr),\qquad
 P_U=p^{-1}(V_U),\qquad Z_U=\pi^{-1}(P_U).
\]
The complement defining \(V_U\) is closed in \(S\), hence in \(V\).
Moreover \(E\cap Z_U=g^{-1}(B_U)\) as topological spaces.
Since the relevant quotients are supported on \(E\), the
affine-section assertion of Proposition~\ref{bl:prop:lifting} gives
\[
 \Gamma(Z_U,I^j/I^d)\longrightarrow
 \Gamma(Z_U,I^j/I^{j+1})
 \quad\text{surjectively}.
\]
Taking \(\mu_r\)-invariants and using
\eqref{bl:lift:invariant-quotients} gives a surjection between sections
over \(P_U\) of the corresponding pulled-back quotient sheaves.
These are canonically linearized pullbacks: their linearizations in
\eqref{bl:lift:invariant-ideals} are the ones obtained by descending
meromorphic functions.

We next take Laurent degree zero for the frame action.  For a coherent
sheaf \(\mathcal A\) on \(V_U\), the affine projection of the
nonzero-vector bundle gives
\begin{equation}\label{bl:lift:laurent-descent}
 p_*p^*\mathcal A
   =\bigoplus_{n\in\Z}\mathcal A\otimes L^{-n}.
\end{equation}
The degree \(n\) term has pullback weight \(rn\).  Sections commute
with this direct sum on the quasi-compact separated open \(V_U\);
this also follows from a finite affine cover.  The preceding
surjection respects the displayed grading.  Its degree-zero map is
therefore surjective, and is exactly
\[
 \Gamma(V_U,J_j^V/J_d^V)\longrightarrow
 \Gamma(V_U,\mathcal B_j).
\]
These are the sections on \(U\) of the two additive pushforwards in
\eqref{bl:lift:descent-surjection}.  Affine opens form a basis, proving
that assertion.

\smallskip\noindent
\emph{Counting the descended layers.}
The sheaves \(Q_j\) are coherent, since \(f\) is projective.
Ceiling arithmetic and the fixed identification \(L|_S=f^*J\) give
\begin{equation}\label{bl:lift:periodicity}
 J_{j-r}^V=J_j^V\otimes L,\qquad
 \mathcal B_{j-r}=\mathcal B_j\otimes f^*J,\qquad
 Q_{j-r}=Q_j\otimes J.
\end{equation}
Also \(J_0^V=\OO_V\) and \(J_1^V=\OO_V(-S)\), so
\(\mathcal B_0=\OO_S\) and \(Q_0=f_*\OO_S\ne0\), because
\(S\ne\varnothing\).

Set
\[
 \mathcal A_N=f_*(J_{-Nr}^V/J_0^V).
\]
By \eqref{bl:lift:descent-surjection} and left exactness, for each
\(-Nr\leq j<0\) there is an exact sequence of additive sheaves
\[
 0\longrightarrow f_*(J_{j+1}^V/J_0^V)
 \longrightarrow f_*(J_j^V/J_0^V)
 \longrightarrow Q_j\longrightarrow0.
\]
Thus \(\mathcal A_N\) has a finite filtration whose coherent
quotients, by \eqref{bl:lift:periodicity}, are precisely
\begin{equation}\label{bl:lift:positive-layers}
 Q_t\otimes J^n,\qquad 0\leq t<r,\quad 1\leq n\leq N.
\end{equation}
The middle additive sheaves need not be coherent \(\OO_T\)-modules;
the following estimates use only these exact sequences and their
coherent quotients.

Choose a common integer \(n_0\geq1\) such that, for all
\(0\leq t<r\) and \(n\geq n_0\), Serre vanishing and generation
give
\[
 H^i(T,Q_t\otimes J^n)=0\ (i>0),\qquad
 Q_t\otimes J^n\text{ is generated by global sections}.
\]
Write \(\ell=\dim T\).  The cohomology sequences of the finite
filtration show that all cohomology of \(\mathcal A_N\) is
finite-dimensional, that it vanishes in degrees greater than \(\ell\),
and, for \(i>0\), that
\begin{equation}\label{bl:lift:higher-bound}
 h^i(T,\mathcal A_N)
 \leq\sum_{t=0}^{r-1}\sum_{n=1}^{n_0-1}
             h^i(T,Q_t\otimes J^n)=:C_i.
\end{equation}
The constants \(C_i\) are independent of \(N\).  Euler
characteristic is additive for the same exact sequences, whence
\[
 \chi(T,\mathcal A_N)
   =\sum_{t=0}^{r-1}\sum_{n=1}^N\chi(T,Q_t\otimes J^n).
\]
For \(n\geq n_0\) all summands are nonnegative, and the summand
with \(t=0\) is at least one: a nonzero globally generated sheaf
has a nonzero global section.  Therefore
\(\chi(T,\mathcal A_N)\to+\infty\).  Together with
\eqref{bl:lift:higher-bound}, this implies
\(h^0(T,\mathcal A_N)\to\infty\).  This argument also covers
zero-dimensional support of the \(Q_t\), and \(\ell=0\): an
individual twist need not grow, since the number of positive layers
already grows with \(N\).

Finally \(J_{-Nr}^V=\OO_V(NG)\), and the definition of the
additive pushforward identifies
\[
 H^0(T,\mathcal A_N)
   =H^0\bigl(V,\OO_V(NG)/\OO_V\bigr).
\]
This proves the claimed growth.
\end{proof}

\begin{proof}[Proof of Proposition~\ref{bl:prop:complex-supported}]
Apply Lemma~\ref{bl:lem:descent-growth}.  From the exact sequence
\[
 0\longrightarrow\OO_V\longrightarrow\OO_V(NG)
   \longrightarrow\OO_V(NG)/\OO_V\longrightarrow0
\]
we obtain
\[
 h^0(V,\OO_V(NG))
 \geq h^0\bigl(V,\OO_V(NG)/\OO_V\bigr)-h^1(V,\OO_V).
\]
The right side is unbounded.  In particular some \(\OO_V(NG)\)
has two linearly independent sections.  On the integral variety \(V\),
their ratio is a nonconstant rational function, so the corresponding
linear system has positive-dimensional image.  Hence
\(\kappa(V,G)>0\).  The actual linear equivalence
\(G\sim qA\) yields \(\kappa(V,A)=\kappa(V,G)>0\), as required.
\end{proof}
